% X_0^6(379)* is hyperelliptic: a Borcherds form on the star, given by its principal part and never
% constructed, makes two CM divisors of degree two linearly equivalent.  Tests in
% tests/DualCertificate6_379_divisor.m and tests/_offline/DualCertificate6_379.m; the solver in
% tools/proof-6-379/.
\documentclass[11pt,reqno]{amsart}
\usepackage[margin=1.15in]{geometry}
\usepackage{amsmath,amssymb,amsthm,mathtools,booktabs}
\usepackage[colorlinks=true,linkcolor=blue,citecolor=blue,urlcolor=blue]{hyperref}
\theoremstyle{plain}
\newtheorem{theorem}{Theorem}[section]
\newtheorem{proposition}[theorem]{Proposition}
\newtheorem{lemma}[theorem]{Lemma}
\theoremstyle{definition}
\newtheorem{remark}[theorem]{Remark}
\DeclareMathOperator{\Jac}{Jac}
\DeclareMathOperator{\Pic}{Pic}
\DeclareMathOperator{\CM}{CM}
\DeclareMathOperator{\Gal}{Gal}
\newcommand{\Q}{\mathbb{Q}}
\newcommand{\Z}{\mathbb{Z}}
\newcommand{\F}{\mathbb{F}}
\begin{document}
\title{$X_0^6(379)^*$ is hyperelliptic}
\author{}
\date{\today}
\maketitle

\section{Statement}

Let $X=X_0^6(379)$ be the Shimura curve attached to an Eichler order of level $379$ in the
quaternion algebra of discriminant $6$, $W=\{w_m : m\mid 2274\}$ its Atkin--Lehner group of order
eight, and $C=X/W$ the star. The curve $C$ has genus three; its Jacobian is simple, of
$\mathrm{GL}_2$-type over a cubic field, and no Atkin--Lehner involution survives on $C$, so none of
the cover arguments applies and hyperellipticity has to be decided on $C$ itself.

\begin{theorem}\label{thm:main}
$C=X_0^6(379)^*$ is hyperelliptic. The hyperelliptic map carries the five fibres
\[
\{P_{-123},P_{-163}\},\quad \{P_{-75},P_{-228}\},\quad \{P_{-52},P_{-168}\},\quad
\{P_{-43},P_{-312}\},\quad \{P_{-147},P_{-408}\},
\]
where $P_d$ is the rational CM point of discriminant $d$ on $C$.
\end{theorem}

Only the first and third fibres are used in the proof; the other three are proved to be fibres by
the same argument (Remark~\ref{rem:others}). Nothing is numerical: the fibres were \emph{found}
by a numerical screen (the canonical images of the two points of a fibre are proportional), and
then \emph{proved} by the certificate below.

\section{CM points on the star}

For a negative discriminant $d$ write $\CM(d)$ for the set of points of $X$ with complex
multiplication by the order $R_d$ of discriminant $d$. It is nonempty exactly when $R_d$ embeds
optimally into the Eichler order, that is, when neither $2$ nor $3$ divides the conductor of $R_d$
or splits in $\Q(\sqrt d)$, and $379$ either divides the conductor or does not stay inert. The
Galois group acts on $\CM(d)$ through $\Pic(R_d)$ by Shimura reciprocity, and the Atkin--Lehner
group acts through orientation changes at the ramified primes together with the classes of the
ideals of norm $m$. Consequently the image $\CM(d)/W$ in $C$ is a single Galois orbit, and its
field of definition has degree one exactly when it is one rational point, which we then denote
$P_d$. The repository computes that degree from $\Pic(R_d)$ and the Atkin--Lehner data alone
(\texttt{DegreeOfFieldOfDefinitionOfCMPoint}); the ten discriminants with a rational point and
class number at most four are
\[
-3,\ -40,\ -43,\ -52,\ -123,\ -163,\ -168,\ -228,\ -312,\ -408,
\]
and $-75=-3\cdot5^2$, $-147=-3\cdot7^2$ are rational as well. None of the four points used below is
fixed by an Atkin--Lehner involution or elliptic, so each has trivial stabiliser.

\section{The certificate}

Let $M=4548=4\cdot1137$, the level of the weight-$\tfrac12$ forms whose Borcherds lifts live on
$X_0^6(379)$ (Guo--Yang \cite{gy}: $M=4\cdot\mathrm{oddpart}(DN)$). A weakly holomorphic form
$f=\sum_nc(n)q^n$ of weight $\tfrac12$ on $\Gamma_0(M)$ with the theta multiplier, holomorphic at
every cusp other than $\infty$, lifts to a meromorphic modular form $\Psi_f$ on $X$, invariant
under $W$ up to a character of finite order, of weight $c(0)$ and with divisor
\begin{equation}\label{eq:div}
\operatorname{div}\Psi_f \;=\; \sum_{n>0}c(-n)\sum_{\substack{r\ge1\\ d=-4n/r^2\ \text{a discriminant}}}
\frac1{e_d}\,P_d
\end{equation}
on $C$, where $P_d$ now stands for the full CM divisor of discriminant $d$ on $C$ (empty if
$\CM(d)=\emptyset$) and $e_d$ is the order of the stabiliser of its points \cite[Lemma~25]{gy}. The
form $f$ is never written down; only its principal part is.

\begin{proposition}\label{prop:cert}
Let $c(-n)$, $1\le n\le700$, be the integers listed in \texttt{data/dual/cert\_6\_379.m} ($208$
of them nonzero). There is a weakly holomorphic form $f$ of weight $\tfrac12$ on $\Gamma_0(4548)$
with the theta multiplier, holomorphic at every cusp other than $\infty$, with principal part
$\sum_nc(-n)q^{-n}$ at $\infty$, and its Borcherds lift is a function on $C$ with divisor
\[
\lambda\,(P_{-123}+P_{-163}-P_{-52}-P_{-168}),\qquad \lambda=308065062 .
\]
\end{proposition}

\begin{proof}
A form of weight $\tfrac12$ with a prescribed principal part at $\infty$ and no other poles exists
if and only if that principal part pairs to zero with every cusp form $g$ of weight $\tfrac32$ on
$\Gamma_0(M)$ with the conjugate multiplier, the multiplier of $\theta^3$: the product $fg\,d\tau$
is then a meromorphic differential on $X_0(M)$ whose only pole is at $\infty$, so its residue
$\sum_nc(-n)a_g(n)$ vanishes, and Serre duality (Riemann--Roch on $X_0(M)$ for the line bundle of
weight-$\tfrac12$ forms) makes the vanishing of all these residues sufficient. That space has
dimension $377$, the obstruction dimension the eta-quotient pipeline predicts at this level, and the
pairing of the listed principal part with a basis of it is zero, exactly, in rational arithmetic
(\texttt{tests/\_offline/DualCertificate6\_379.m}; the test also checks that the pairing is not
identically zero, by perturbing one coefficient). The divisor of the lift is \eqref{eq:div}: the
listed coefficients sum to $\lambda$ at $d=-123$ and $-163$, to $-\lambda$ at $-52$ and $-168$, and
to $0$ at each of the other $21$ discriminants carrying CM points that are reachable from poles of
order at most $700$ (\texttt{tests/DualCertificate6\_379\_divisor.m}); the four stabilisers are
trivial. The divisor has degree zero on $C$, while a Borcherds form of weight $k$ has a divisor of
degree $k$ times the hyperbolic area of $C$ over $4\pi$, so the lift has weight zero and is a
function up to a character of finite order; some power of it is a genuine function on $C$.
\end{proof}

How the vector was found: the unknowns are the $c(-n)$, $n\le700$; the equations are the
$377$ pairings and, for each of the $25$ discriminants $d$ reachable from poles of order at most
$700$, the condition that the coefficients sum to $e_d$ times the prescribed multiplicity. The
system is consistent, a rational solution was cleared of denominators and made primitive
(\texttt{tools/proof-6-379/stepH-dual.m}). It is not consistent with poles of order at most $500$,
and it is not consistent for the divisor $P_{-123}-P_{-163}$, two points of one fibre, as it must
not be: that class is a nonzero point of a torsion-free group.

\section{The torsion, and the proof of the theorem}

\begin{lemma}\label{lem:tors}
$\Jac(C)(\Q)$ is torsion-free.
\end{lemma}

\begin{proof}
For an odd prime $p$ of good reduction, reduction is injective on $\Jac(C)(\Q)_{\mathrm{tors}}$,
so the torsion order divides $\#\Jac(C)(\F_p)$. The latter is $P(1)$ for the characteristic
polynomial $P$ of Frobenius, which the trace formula gives from the point counts of $C$ over
$\F_p$, $\F_{p^2}$, $\F_{p^3}$ (Newton's identities and the functional equation), with no model of
$C$. The greatest common divisor over the good primes $5\le p\le59$ is $1$
(\texttt{tests/DualCertificate6\_379\_divisor.m}).
\end{proof}

\begin{proof}[Proof of Theorem~\ref{thm:main}]
By Proposition~\ref{prop:cert} the class of $\Delta=P_{-123}+P_{-163}-P_{-52}-P_{-168}$ in
$\Jac(C)(\Q)$ is killed by a positive integer, so it is torsion, so it is zero by
Lemma~\ref{lem:tors}. Thus $P_{-123}+P_{-163}$ and $P_{-52}+P_{-168}$ are two distinct effective
divisors of degree two in one linear equivalence class, that class has $h^0\ge2$, and $C$ carries
a $g^1_2$: it is hyperelliptic, and both divisors are fibres of the hyperelliptic map.
\end{proof}

\begin{remark}\label{rem:others}
The same linear system is consistent for $P_{-75}+P_{-228}$, $P_{-43}+P_{-312}$ and
$P_{-147}+P_{-408}$ against $P_{-52}+P_{-168}$ (poles of order at most $2400$), so all five
pairs of Theorem~\ref{thm:main} are fibres; these three certificates are not banked.
\end{remark}

\begin{remark}
Nothing here produces an equation for $C$. The proof is of the shape of \texttt{paper/torsion-counting-hyperelliptic.tex} (the star of $X_0^{46}(31)$): a relation between CM divisors, exact up to torsion, against a torsion bound. What is new is how the relation is certified. The eta-quotient search for a Borcherds form with this divisor ran for a day and a half at level $4548$ without finishing; the dual formulation decides the existence of the form by linear algebra against the weight-$\tfrac32$ cusp forms, in four hours for the basis and seconds for the solve, and the same basis then decides every other candidate divisor in seconds.
\end{remark}

\begin{thebibliography}{9}
\bibitem{gy} J.~Guo and Y.~Yang, \emph{Equations of hyperelliptic Shimura curves}, Compos. Math.
153 (2017), 1–40.
\end{thebibliography}
\end{document}
